%==============================================================
%  lecons/ensembles/construction_C.tex — Construction de C
%==============================================================

\lecontitle{Construction de \texorpdfstring{$\mathbb{C}$}{C}}{Théorie des ensembles — Clôture algébrique de \texorpdfstring{$\mathbb{R}$}{R}}

%=============================================================
%  § 1 — MOTIVATION
%=============================================================
\secnum{navyblue}{1}{Motivation — fermer algébriquement $\mathbb{R}$}

\rubriq{Le problème}
Dans $\mathbb{R}$, le polynôme $X^2 + 1$ n'a pas de racine. On veut étendre
$\mathbb{R}$ en un corps algébriquement clos, c'est-à-dire dans lequel
tout polynôme non constant a au moins une racine.

\begin{tcolorbox}[thmbox]
  \textbf{Construction algébrique.}\enspace%
  \index{nombres complexes!construction}\index{corps algébriquement clos}
  On construit $\mathbb{C}\cong\mathbb{R}[X]/(X^2+1)$ en réalisant cet isomorphisme
  de corps sur l'ensemble $\mathbb{R}^2$ (à ne pas confondre avec l'anneau-produit
  $\mathbb{R}\times\mathbb{R}$ qui n'est pas un corps), muni des lois :
  \begin{align*}
    (a,b) + (c,d) &= (a+c,\; b+d), \\
    (a,b) \times (c,d) &= (ac-bd,\; ad+bc).
  \end{align*}
  On note $i = (0,1)$, de sorte que $i^2 = (-1, 0)$, identifié à $-1$.
  L'injection $\mathbb{R} \hookrightarrow \mathbb{C},\; a \mapsto (a,0)$,
  est un homomorphisme de corps.
\end{tcolorbox}

\begin{significationintuitive}
  Le quotient $\mathbb{R}[X]/(X^2+1)$ consiste à manipuler les polynômes réels
  \emph{modulo} la relation $X^2+1 = 0$, c'est-à-dire en convenant que $X^2 = -1$.
  Tout reste de division par $X^2+1$ est de degré $\leq 1$, donc de la forme
  $a + bX$ : on retrouve ainsi le plan $\mathbb{R}^2$, et la classe de $X$ joue le
  rôle de $i$. La table de multiplication $(a,b)\times(c,d) = (ac-bd,\,ad+bc)$
  n'est rien d'autre que le produit $(a+bX)(c+dX) = ac + (ad+bc)X + bd\,X^2$ où
  l'on remplace $X^2$ par $-1$.
\end{significationintuitive}

\decsep{}

%=============================================================
%  § 2 — STRUCTURE ALGÉBRIQUE
%=============================================================
\secnum{navyblue}{2}{$\mathbb{C}$ est un corps}

\begin{mdframed}[style=proofbar]
  \textit{Corps.} Les axiomes de groupe abélien pour $+$ sont immédiats. La
  multiplication est commutative, associative et distributive sur $+$ : cela se
  vérifie par calcul direct à partir des propriétés correspondantes de $\mathbb{R}$,
  l'élément neutre étant $(1,0)$. Pour l'inverse multiplicatif de
  $(a,b) \neq (0,0)$ :
  \[
    {(a,b)}^{-1} = \left(\frac{a}{a^2+b^2},\; \frac{-b}{a^2+b^2}\right).
  \]
  En effet $(a,b)\cdot(a,-b) = (a^2+b^2, 0)$ s'identifie au réel $a^2+b^2 \neq 0$
  (car $a^2+b^2 > 0$ dès que $(a,b)\neq(0,0)$), d'où la formule ci-dessus en
  divisant par ce réel.

  \textit{Dimension.} $\mathbb{C}$ est un $\mathbb{R}$-espace vectoriel de dimension $2$,
  de base $(1, i)$. Tout $z \in \mathbb{C}$ s'écrit $z = a + ib$ avec $a,b\in\mathbb{R}$.
  $\blacksquare$
\end{mdframed}

\begin{tcolorbox}[thmbox]
  \textbf{Propriétés fondamentales.}\enspace%
  \index{nombres complexes!module}\index{nombres complexes!conjugué}
  Pour $z = a+ib$, on définit :
  \[
    \bar{z} = a - ib \quad \text{(conjugué)}, \qquad
    |z| = \sqrt{a^2+b^2} \quad \text{(module)}.
  \]
  On a $z \bar{z} = |z|^2$, $|zw| = |z||w|$, et $|z+w| \leq |z|+|w|$
  (inégalité triangulaire).
  La forme polaire est $z = |z|e^{i\theta}$ avec $\theta = \arg z$.
\end{tcolorbox}

\decsep{}

%=============================================================
%  § 3 — THÉORÈME FONDAMENTAL DE L'ALGÈBRE
%=============================================================
\secnum{navyblue}{3}{Théorème fondamental de l'algèbre}

\begin{tcolorbox}[thmbox]
  \textbf{Théorème (d'Alembert–Gauss).}\enspace%
  \index{théorème fondamental de l'algèbre}
  $\mathbb{C}$ est \emph{algébriquement clos} : tout polynôme
  $P \in \mathbb{C}[X]$ de degré $n \geq 1$ admet au moins une racine dans
  $\mathbb{C}$.

  \medskip
  \emph{Corollaire.} Par récurrence sur le degré (en factorisant un facteur
  $(X-a)$ à chaque étape), tout $P$ de degré $n \geq 1$ admet alors
  \emph{exactement} $n$ racines dans $\mathbb{C}$, comptées avec multiplicités.

  \medskip
  Comme $\mathbb{C}$ est algébriquement clos et que l'extension
  $\mathbb{C}/\mathbb{R}$ est algébrique (de degré $2$), $\mathbb{C}$ est la
  \emph{clôture algébrique} de $\mathbb{R}$. À ne pas confondre avec la définition
  d'algébriquement clos elle-même, qui exprime que toute extension algébrique de
  $\mathbb{C}$ est triviale (égale à $\mathbb{C}$).
\end{tcolorbox}

\begin{mdframed}[style=proofbar]
  \textit{Esquisse (preuve topologique).}
  Quitte à diviser $P$ par son coefficient dominant (ce qui ne change pas ses
  racines), on suppose $P$ \emph{unitaire}. Supposons que $P$ (de degré
  $n\geq 1$) n'ait aucune racine. Alors
  $z \mapsto P(z)/|P(z)|$ est définie et continue sur $\mathbb{C}$ tout entier,
  à valeurs dans le cercle $\mathbb{S}^1$. Sa restriction au cercle de rayon $R$
  est donc une boucle, dont l'indice (élément de $\pi_1(\mathbb{S}^1)\cong\mathbb{Z}$)
  ne dépend pas de $R$~— la boucle varie continûment avec $R$ et son indice,
  entier, est donc constant : en faisant $R \to 0$, la boucle se contracte
  en la constante $P(0)/|P(0)|$, donc cet indice est nul. Mais pour $|z|\to\infty$
  on a $P(z)/|P(z)| \sim z^n/|z|^n$, qui s'enroule $n$ fois autour de l'origine,
  d'indice $n$. Comme l'indice est un entier invariant par homotopie, on aurait
  $0 = n$ : contradiction puisque $n \geq 1$. \hfill$\blacksquare$
\end{mdframed}

\rubriq{Conséquence : classification des corps}
Toute extension algébrique finie de $\mathbb{R}$ est soit $\mathbb{R}$ soit $\mathbb{C}$
(extension de degré $2$). Il n'existe pas d'extension algébrique de $\mathbb{R}$
de degré $3$ ou plus.

\decsep{}

%=============================================================
%  § 4 — EXEMPLES, PIÈGES ET EXERCICES
%=============================================================
\secnum{exgreen}{4}{Exemples, pièges et exercices}

\rubriq{Exemples}

\begin{mdframed}[style=exbar]
  \textbf{1. Formule d'Euler.}\enspace%
  \index{formule d'Euler}
  $e^{i\theta} = \cos\theta + i\sin\theta$. En particulier,
  $e^{i\pi} + 1 = 0$ (identité d'Euler).

  \smallskip\noindent
  \textbf{2. Racines $n$-ièmes de l'unité.}\enspace%
  \index{racines de l'unité}
  Les solutions de $z^n = 1$ sont $\omega_k = e^{2ik\pi/n}$ pour $k = 0, \ldots, n-1$.
  Elles forment le groupe cyclique $\mu_n \cong \mathbb{Z}/n\mathbb{Z}$.

  \smallskip\noindent
  \textbf{3. Transformée de Fourier.}\enspace%
  L'utilisation de $\mathbb{C}$ est fondamentale en analyse harmonique :
  $\hat{f}(\xi) = \int_{\mathbb{R}} f(x) e^{-2\pi i x\xi}\,\mathrm{d}x$.
\end{mdframed}

\vspace{6pt}
\rubriq{Pièges}

\begin{mdframed}[style=cexbar]
  \textbf{$\mathbb{C}$ n'est pas un corps ordonné.}\enspace%
  Il n'existe aucun ordre total sur $\mathbb{C}$ compatible avec $+$ et $\times$.
  En effet, dans un corps ordonné tout carré est $\geq 0$, donc on aurait
  $i^2 = -1 \geq 0$, en contradiction avec $-1 < 0$.

  \smallskip\noindent
  \textbf{Pas de clôture algébrique canonique.}\enspace%
  La clôture algébrique de $\mathbb{Q}$ (les nombres algébriques $\bar{\mathbb{Q}}$)
  n'est pas $\mathbb{C}$ entier — c'est un sous-corps strict. Par ailleurs, la
  clôture algébrique est unique \emph{à isomorphisme près}, mais cet isomorphisme
  n'est ni unique ni canonique : le groupe $\mathrm{Aut}(\bar{\mathbb{Q}}/\mathbb{Q})$
  est immense.
\end{mdframed}

\vspace{6pt}
\exolabel

\begin{mdframed}[style=exobar]
  \begin{enumerate}
    \item Montrer que $\mathbb{R}[X]/(X^2+1)$ est bien un corps, et que
          $(X^2+1)$ est un idéal maximal de $\mathbb{R}[X]$.

    \item Trouver toutes les racines carrées de $3 + 4i$.

    \item Montrer que tout automorphisme de corps de $\mathbb{C}$
          fixant $\mathbb{R}$ est soit l'identité, soit la conjugaison.

    \item Factoriser $X^4 + 4$ dans $\mathbb{R}[X]$, puis dans $\mathbb{C}[X]$.

    \item (Théorème de Liouville) Montrer qu'une fonction entière bornée
          est constante, et en déduire le théorème fondamental de l'algèbre.
  \end{enumerate}
\end{mdframed}

\leconfooter{R.\ Descartes (1637) — C.\ F.\ Gauss (1799, thèse) — J.\ le Rond d'Alembert (1746)}
